\begin{tikzpicture}[font=\small,>=stealth]
\draw[gray!45,rounded corners=3pt] (0,0) rectangle (14.2,9.5);
\node at (7.1,8.85) {先不画双曲线：同一直线上的长度比，等于水平投影比};

\begin{scope}[shift={(1.7,2.2)},scale=.85]
\coordinate (A) at (0,0);\coordinate (B) at (3,4);\coordinate (C) at (4.5,6);
\fill[blue!12] (A)--(3,0)--(B)--cycle;
\fill[orange!18] (B)--(4.5,4)--(C)--cycle;
\draw[blue!75!black,thick] (A)--(3,0)--(B)--cycle;
\draw[orange!85!black,thick] (B)--(4.5,4)--(C)--cycle;
\draw[gray] (2.8,0)--(2.8,.2)--(3,.2) (4.3,4)--(4.3,4.2)--(4.5,4.2);
\fill (A) circle(2.8pt) node[left=5pt] {$A$};
\fill (B) circle(2.8pt) node[above left=5pt] {$B$};
\fill (C) circle(2.8pt) node[right=5pt] {$C$};
\node[blue!75!black,left=5pt] at (1.5,2) {$5$};
\node[orange!85!black,left=5pt] at (3.75,5) {$5/2$};
\node[blue!75!black,below=5pt] at (1.5,0) {$3$};
\node[orange!85!black,below=5pt] at (3.75,4) {$3/2$};
\end{scope}
\node[anchor=west,align=left] at (8,6.5) {两个直角三角形相似\\[14pt]$\displaystyle\frac{AB}{BC}=\frac5{5/2}=2$\\[14pt]$\displaystyle\frac{\text{横向差}}{\text{横向差}}=\frac3{3/2}=2$};
\node[anchor=west,align=left] at (8,3.1) {同一方向、同一放大倍数\\[10pt]取比时，公共倍数约掉。};
\node at (7.1,.65) {接下来把蓝、橙两段对应为 $FM,MN$；原理不变。};

\end{tikzpicture}
